\section{Regular field evolution and energy}\label{sec:energy}

The prepared amplitude in \eqref{eq:realization} has two convenient
representations: the positive scalar momentum envelope \(F_\lambda(p)\)
and the real scalar position envelope \(\sqrt{2/\pi}\,h(\tau,x)\),
each multiplying the fixed angular spinor \(\chi\).
Positivity in the first representation does not imply positivity in the
second. This distinction is already enough to make nodal events
compatible with a nonnegative generator.

\subsection{Global mass and kinetic energy}
Define
\[
 N(\tau)=\norm{\Psi_\tau}^2,\qquad
 T(\tau)=\inner{\Psi_\tau}{K\Psi_\tau}.
\]
The fixed finite-order bounds of Section~\ref{sec:source} justify every
differentiation and integration by parts in this section.

\begin{proposition}[Global dissipation]\label{prop:global-energy}
The complete prepared orbit satisfies
\begin{align}
 N&=\int_\R |F_{\lambda_0-\tau}(p)|^2\dd p
   =\frac2\pi\int_\R h^2\dd x>0,\label{eq:N}\\
 T&=\int_\R p^2|F_{\lambda_0-\tau}(p)|^2\dd p
   =\frac2\pi\int_\R h_x^2\dd x,\label{eq:T}\\
 N'&=-2T,\qquad T'=-2\norm{K\Psi_\tau}^2.\label{eq:dissipation}
\end{align}
For the normalized field \(\phi=\Psi/\sqrt N\) and mean relative
energy \(E=\inner\phi{K\phi}\),
\begin{equation}\label{eq:normalized}
 \phi'=-(K-E)\phi,\qquad
 E'=-2\Var_\phi(K),\qquad
 \norm{\phi'}^2=\Var_\phi(K),
\end{equation}
where \(\Var_\phi(K)=\inner\phi{K^2\phi}-E^2\).
All these quantities are finite and smooth on compact time intervals.
\end{proposition}
\begin{proof}
Plancherel and the first-shell multiplier give \eqref{eq:N}--\eqref{eq:T}.
Differentiating \(\Psi'=-K\Psi\), and using self-adjointness, proves
\eqref{eq:dissipation}. The source is nonzero, so \(N>0\).
The quotient rule gives
\(\phi'=-K\phi-(N'/2N)\phi=-(K-E)\phi\).
Differentiating \(E=T/N\) gives
\[
 E'=-2\frac{\norm{K\Psi}^2}{N}+2\frac{T^2}{N^2}
    =-2\Var_\phi(K).
\]
Taking the squared norm of the first identity proves the last one.
All products are legitimate because the prepared field lies in every
fixed power domain. The source bounds also give their local uniformity.
\end{proof}

The probability measure used in these spectral expectations is
\begin{equation}\label{eq:spectral-measure}
 \dd\nu_\tau(p)=\frac{|F_{\lambda_0-\tau}(p)|^2}{N(\tau)}\dd p.
\end{equation}
Its mean \(p^2\) is \(E\). It differs from the linear source measure
\(F_\lambda(p)\dd p/\int F_\lambda\), which will describe the cosine
balances. In particular the Gaussian exponent in \(\nu_\tau\) is twice
the exponent in that linear measure. Confusing the two loses a factor
of two in the evolution and changes the quantity being estimated.

\subsection{Local intensity balance}
Since \(h\) is real for real \((\tau,x)\),
\begin{equation}\label{eq:local-density}
 \partial_\tau(h^2)=\partial_x^2(h^2)-2h_x^2.
\end{equation}
Indeed, \(2hh_\tau=2hh_{xx}=(h^2)_{xx}-2h_x^2\).
The corresponding angular-integrated field intensity is
\((2/\pi)h^2\). Thus local intensity redistributes and dissipates under
imaginary time; its integral is not a conserved unitary probability.

At a double node \(h=h_x=0\), \(h_{xx}=a\ne0\), both
\(\partial_\tau(h^2)\) and \(\partial_x^2(h^2)\) vanish, while
\begin{equation}\label{eq:density-event}
 \partial_\tau^2(h^2)=2a^2.
\end{equation}
The field itself has nonzero time derivative \(h_\tau=a\).
This is the first indication that a smooth field can pass through a
degenerate zero without producing an intensity or global-energy spike.
The spatial and local-integral scales are derived below.
